\begin{frame}{2022$\sim$2023 --- 정렬(alignment)과 프론티어 모델}
  \small
  \begin{itemize}
    \item \kb{InstructGPT (Ouyang et al., 2022)} --- \textbf{RLHF}(인간
      피드백 강화학습)로 언어모델을 ``사람이 원하는 방식''에 맞춤 ---
      \textbf{ChatGPT를 실제로 가능하게 한 기법}.
    \item \kb{GPT-4 Technical Report (OpenAI, 2023)} --- 대규모
      멀티모달 모델. ``\textbf{프론티어 모델}(frontier model)''이라는
      용어가 이 시기부터 자리잡음.
    \item \kb{DALL-E 2 (Ramesh et al., 2022)} --- CLIP + 디퓨전으로
      텍스트 $\to$ 이미지 생성. ChatGPT와는 \textbf{다른 갈래}(생성모델,
      Week 15)지만 같은 시대의 돌파구.
  \end{itemize}
  \vskip0.5em
  \kq{공통점: 근본적으로 새 수학이 아니라, Block A/B에서 배울
  \textbf{기존 아이디어(어텐션, 강화학습, 디퓨전)를 더 큰 규모로
  조합}한 것.}
  \vskip0.4em
  {\tiny\color{ksagray} Ouyang, L. et al. (2022). Training language
  models to follow instructions with human feedback. \textit{NeurIPS}.
  \\ OpenAI (2023). GPT-4 Technical Report. \textit{arXiv:2303.08774}.
  \\ Ramesh, A. et al. (2022). Hierarchical Text-Conditional Image
  Generation with CLIP Latents. \textit{arXiv:2204.06125}.}
\end{frame}
